\documentclass[11pt,a4paper]{article}
\usepackage[utf8]{inputenc}
\usepackage[T1]{fontenc}
\usepackage[catalan]{babel}
\usepackage{lmodern}
\usepackage[sfdefault,lining]{FiraSans}
\usepackage[a4paper,top=2.3cm,bottom=2.2cm,left=1.8cm,right=1.8cm,headheight=1.3cm,headsep=0.4cm,footskip=1cm]{geometry}
\usepackage{amsmath,amssymb}
\usepackage{icomma}
\usepackage{graphicx}
\usepackage[table]{xcolor}
\usepackage{tikz}
\usepackage{fontawesome5}
\usepackage{ccicons}
\usepackage{fancyhdr}
\usepackage{lastpage}
\usepackage{tasks}
\usepackage{needspace}
\usepackage[most]{tcolorbox}
\usepackage{array}
\usepackage[hidelinks,pdftitle={Fitxa d'exercicis de la unitat 6: Sistemes d'equacions. 3r d'ESO},pdfsubject={Matemàtiques 3r ESO},pdfkeywords={sistemes d'equacions, equacions lineals amb dues incògnites, substitució, igualació, reducció, mètode gràfic, classificació de sistemes, problemes}]{hyperref}

\definecolor{navy}{HTML}{1B2748}
\definecolor{accent}{HTML}{E8603A}
\definecolor{teal}{HTML}{1F9E90}
\definecolor{soft}{HTML}{EEF1F8}
\definecolor{muted}{HTML}{5D6685}

\setlength{\parindent}{0pt}
\setlength{\parskip}{3pt}
\settasks{label=\alph*), label-format=\bfseries\color{navy}, label-width=1.3em, item-indent=1.8em, column-sep=0.6em, after-item-skip=5pt, before-skip=3pt, after-skip=4pt}

% ---------- capçalera i peu ----------
\newcommand{\llicencia}{\href{https://creativecommons.org/licenses/by-nc-sa/4.0/deed.ca}{\ccbyncsa}\ \ Llicència Creative Commons Reconeixement-NoComercial-CompartirIgual 4.0 Internacional (CC BY-NC-SA 4.0)}
\pagestyle{fancy}
\fancyhf{}
\fancyhead[L]{\raisebox{-0.35\height}{\includegraphics[height=0.95cm]{logo-gva.png}}}
\fancyhead[C]{\small\color{muted}Fitxa d'exercicis · Unitat 6: Sistemes d'equacions · 3r d'ESO}
\fancyhead[R]{\raisebox{-0.35\height}{\includegraphics[height=0.95cm]{logo-vermellar.png}}}
\fancyfoot[L]{\scriptsize\color{muted}\llicencia}
\fancyfoot[R]{\small\color{muted}\thepage}
\renewcommand{\headrulewidth}{0.4pt}
\renewcommand{\headrule}{\color{muted!40}\hrule height \headrulewidth}
\fancypagestyle{primera}{\fancyhf{}\renewcommand{\headrulewidth}{0pt}
  \fancyfoot[L]{\scriptsize\color{muted}\llicencia}
  \fancyfoot[R]{\small\color{muted}\thepage}}

% ---------- nivells i exercicis ----------
\newcommand{\estrelles}[1]{{\color{accent}\ifcase#1\or\faStar\faStar[regular]\faStar[regular]\or\faStar\faStar\faStar[regular]\or\faStar\faStar\faStar\fi}}
\newcounter{exer}
\newcommand{\ex}[1]{\par\needspace{4\baselineskip}\medskip\stepcounter{exer}\noindent{\color{navy}\bfseries\theexer.}\ \ {\small\estrelles{#1}}\ \ }

\newcommand{\seccio}[2]{\par\needspace{10\baselineskip}\bigskip
  \begin{tcolorbox}[colback=navy,colframe=navy,coltext=white,arc=2mm,boxrule=0pt,left=3mm,right=3mm,top=1.2mm,bottom=1.2mm]
  {\large\bfseries #1\quad #2}
  \end{tcolorbox}\vspace{-2pt}}

\newcolumntype{C}[1]{>{\centering\arraybackslash}p{#1}}
\usepackage{pifont}

% quadre de recordatori (només a la fitxa d'exercicis)
\newcommand{\recorda}[1]{\par\smallskip\begin{tcolorbox}[colback=accent!7,colframe=accent!55,arc=2mm,boxrule=0.5pt,left=2.5mm,right=2.5mm,top=1mm,bottom=1mm]\small{\color{accent}\faLightbulb[regular]}\ \textbf{Recorda:}\ #1\end{tcolorbox}}
\newcommand{\res}[1]{#1}
\newcommand{\rest}[1]{#1}

\def\fmneed{14}
% =====================================================================
% Capçalera FentMates de les fitxes: un color per a cada curs.
% Els colors són els de la pàgina inicial (LEVELS a assets/js/materials.js);
% el logo ve de _eines/logo/pdf/ (python3 _eines/logo/genera.py).
%
% Ús, al final del preàmbul de la unitat:
%   % =====================================================================
% Capçalera FentMates de les fitxes: un color per a cada curs.
% Els colors són els de la pàgina inicial (LEVELS a assets/js/materials.js);
% el logo ve de _eines/logo/pdf/ (python3 _eines/logo/genera.py).
%
% Ús, al final del preàmbul de la unitat:
%   % =====================================================================
% Capçalera FentMates de les fitxes: un color per a cada curs.
% Els colors són els de la pàgina inicial (LEVELS a assets/js/materials.js);
% el logo ve de _eines/logo/pdf/ (python3 _eines/logo/genera.py).
%
% Ús, al final del preàmbul de la unitat:
%   \input{../comu/fentmates.tex}
%   \fmcurs{1eso}
% i, al començament del document:
%   \capcalera{Matemàtiques}{1r ESO}{Fitxa d'exercicis · Unitat 2}{Divisibilitat}
% Opcional abans de \input: \def\fmneed{14} (espai mínim abans de cada secció).
% =====================================================================
\makeatletter
\def\fm@logodir{../../logo/pdf}
\providecommand{\fmneed}{10}
\def\fm@cursnom{}

% clau del curs, color i color del text damunt del color
\newcommand{\fm@curs}[3]{\@namedef{fm@color@#1}{#2}\@namedef{fm@text@#1}{#3}}
\fm@curs{1eso}{E6B84A}{navy}
\fm@curs{2eso}{3B82F6}{white}
\fm@curs{3eso}{FF7A50}{white}
\fm@curs{4eso-a}{A78BFA}{white}
\fm@curs{4eso-b}{6D28D9}{white}
\fm@curs{1batx-mi}{9B1C1C}{white}
\fm@curs{1batx-ccss}{EF6B6B}{white}
\fm@curs{2batx-mii}{1E6B47}{white}
\fm@curs{2batx-ccss}{2BB3A3}{white}

\newcommand{\fmcurs}[1]{%
  \@ifundefined{fm@color@#1}{\PackageError{fentmates}{Curs desconegut: #1}{}}{}%
  \edef\fm@tmp{\noexpand\definecolor{curs}{HTML}{\@nameuse{fm@color@#1}}%
               \noexpand\colorlet{curstext}{\@nameuse{fm@text@#1}}}\fm@tmp
  \colorlet{cursfosc}{curs!72!black}%
  \def\fm@logo{\fm@logodir/vertical-#1.pdf}%
  \def\fm@marca{\fm@logodir/marca-#1.pdf}}

% ---------- capçalera de la primera pàgina ----------
% #1 matèria (petit), #2 curs, #3 tipus de document i unitat, #4 títol
\newcommand{\capcalera}[4]{%
  \gdef\fm@cursnom{#2}%
  \par\noindent
  \begin{tcolorbox}[colback=curs,colframe=curs,boxrule=0pt,arc=2.2mm,width=3.5cm,height=2.15cm,
      valign=center,left=3mm,right=2mm,top=1mm,bottom=1mm,nobeforeafter,box align=center]
    \color{curstext}{\footnotesize #1}\par\vspace{1pt}{\fontsize{19}{21}\selectfont\bfseries #2\par}%
  \end{tcolorbox}\hspace{5mm}%
  \begin{minipage}[c]{\dimexpr\linewidth-3.5cm-5mm-2.9cm\relax}
    {\color{cursfosc}\small\bfseries\MakeUppercase{#3}}\par\vspace{2pt}
    {\color{navy}\fontsize{20}{23}\selectfont\bfseries #4\par}
  \end{minipage}\hfill
  \begin{minipage}[c]{2.6cm}\raggedleft\includegraphics[width=2.6cm]{\fm@logo}\end{minipage}\par
  \vspace{2pt}}

% ---------- barres de secció del color del curs ----------
\renewcommand{\seccio}[2]{\par\needspace{\fmneed\baselineskip}\bigskip
  \begin{tcolorbox}[colback=curs,colframe=curs,coltext=curstext,arc=2mm,boxrule=0pt,left=3mm,right=3mm,top=1.2mm,bottom=1.2mm]
  {\large\bfseries #1\quad #2}
  \end{tcolorbox}\vspace{-2pt}}

% ---------- capçalera i peu de les pàgines ----------
% primera pàgina: logos institucionals al peu
\fancypagestyle{primera}{\fancyhf{}\renewcommand{\headrulewidth}{0pt}
  \fancyfoot[L]{\raisebox{-0.32\height}{\includegraphics[height=0.78cm]{logo-gva.png}}\hspace{5mm}%
    \raisebox{-0.32\height}{\includegraphics[height=0.78cm]{logo-vermellar.png}}}
  \fancyfoot[R]{\scriptsize\color{muted}\href{https://creativecommons.org/licenses/by-nc-sa/4.0/deed.ca}{\ccbyncsa}\ \ CC BY-NC-SA 4.0 · fentmates.org\qquad\small\thepage}}
% la resta: el símbol a l'esquerra i el curs a la dreta, amb el títol al centre com abans
\fancyhead[L]{\raisebox{-0.3\height}{\includegraphics[height=0.62cm]{\fm@marca}}}
\fancyhead[R]{\small\bfseries\color{cursfosc}\fm@cursnom}
\renewcommand{\headrule}{\color{curs!45}\hrule height \headrulewidth}
\makeatother

%   \fmcurs{1eso}
% i, al començament del document:
%   \capcalera{Matemàtiques}{1r ESO}{Fitxa d'exercicis · Unitat 2}{Divisibilitat}
% Opcional abans de \input: \def\fmneed{14} (espai mínim abans de cada secció).
% =====================================================================
\makeatletter
\def\fm@logodir{../../logo/pdf}
\providecommand{\fmneed}{10}
\def\fm@cursnom{}

% clau del curs, color i color del text damunt del color
\newcommand{\fm@curs}[3]{\@namedef{fm@color@#1}{#2}\@namedef{fm@text@#1}{#3}}
\fm@curs{1eso}{E6B84A}{navy}
\fm@curs{2eso}{3B82F6}{white}
\fm@curs{3eso}{FF7A50}{white}
\fm@curs{4eso-a}{A78BFA}{white}
\fm@curs{4eso-b}{6D28D9}{white}
\fm@curs{1batx-mi}{9B1C1C}{white}
\fm@curs{1batx-ccss}{EF6B6B}{white}
\fm@curs{2batx-mii}{1E6B47}{white}
\fm@curs{2batx-ccss}{2BB3A3}{white}

\newcommand{\fmcurs}[1]{%
  \@ifundefined{fm@color@#1}{\PackageError{fentmates}{Curs desconegut: #1}{}}{}%
  \edef\fm@tmp{\noexpand\definecolor{curs}{HTML}{\@nameuse{fm@color@#1}}%
               \noexpand\colorlet{curstext}{\@nameuse{fm@text@#1}}}\fm@tmp
  \colorlet{cursfosc}{curs!72!black}%
  \def\fm@logo{\fm@logodir/vertical-#1.pdf}%
  \def\fm@marca{\fm@logodir/marca-#1.pdf}}

% ---------- capçalera de la primera pàgina ----------
% #1 matèria (petit), #2 curs, #3 tipus de document i unitat, #4 títol
\newcommand{\capcalera}[4]{%
  \gdef\fm@cursnom{#2}%
  \par\noindent
  \begin{tcolorbox}[colback=curs,colframe=curs,boxrule=0pt,arc=2.2mm,width=3.5cm,height=2.15cm,
      valign=center,left=3mm,right=2mm,top=1mm,bottom=1mm,nobeforeafter,box align=center]
    \color{curstext}{\footnotesize #1}\par\vspace{1pt}{\fontsize{19}{21}\selectfont\bfseries #2\par}%
  \end{tcolorbox}\hspace{5mm}%
  \begin{minipage}[c]{\dimexpr\linewidth-3.5cm-5mm-2.9cm\relax}
    {\color{cursfosc}\small\bfseries\MakeUppercase{#3}}\par\vspace{2pt}
    {\color{navy}\fontsize{20}{23}\selectfont\bfseries #4\par}
  \end{minipage}\hfill
  \begin{minipage}[c]{2.6cm}\raggedleft\includegraphics[width=2.6cm]{\fm@logo}\end{minipage}\par
  \vspace{2pt}}

% ---------- barres de secció del color del curs ----------
\renewcommand{\seccio}[2]{\par\needspace{\fmneed\baselineskip}\bigskip
  \begin{tcolorbox}[colback=curs,colframe=curs,coltext=curstext,arc=2mm,boxrule=0pt,left=3mm,right=3mm,top=1.2mm,bottom=1.2mm]
  {\large\bfseries #1\quad #2}
  \end{tcolorbox}\vspace{-2pt}}

% ---------- capçalera i peu de les pàgines ----------
% primera pàgina: logos institucionals al peu
\fancypagestyle{primera}{\fancyhf{}\renewcommand{\headrulewidth}{0pt}
  \fancyfoot[L]{\raisebox{-0.32\height}{\includegraphics[height=0.78cm]{logo-gva.png}}\hspace{5mm}%
    \raisebox{-0.32\height}{\includegraphics[height=0.78cm]{logo-vermellar.png}}}
  \fancyfoot[R]{\scriptsize\color{muted}\href{https://creativecommons.org/licenses/by-nc-sa/4.0/deed.ca}{\ccbyncsa}\ \ CC BY-NC-SA 4.0 · fentmates.org\qquad\small\thepage}}
% la resta: el símbol a l'esquerra i el curs a la dreta, amb el títol al centre com abans
\fancyhead[L]{\raisebox{-0.3\height}{\includegraphics[height=0.62cm]{\fm@marca}}}
\fancyhead[R]{\small\bfseries\color{cursfosc}\fm@cursnom}
\renewcommand{\headrule}{\color{curs!45}\hrule height \headrulewidth}
\makeatother

%   \fmcurs{1eso}
% i, al començament del document:
%   \capcalera{Matemàtiques}{1r ESO}{Fitxa d'exercicis · Unitat 2}{Divisibilitat}
% Opcional abans de \input: \def\fmneed{14} (espai mínim abans de cada secció).
% =====================================================================
\makeatletter
\def\fm@logodir{../../logo/pdf}
\providecommand{\fmneed}{10}
\def\fm@cursnom{}

% clau del curs, color i color del text damunt del color
\newcommand{\fm@curs}[3]{\@namedef{fm@color@#1}{#2}\@namedef{fm@text@#1}{#3}}
\fm@curs{1eso}{E6B84A}{navy}
\fm@curs{2eso}{3B82F6}{white}
\fm@curs{3eso}{FF7A50}{white}
\fm@curs{4eso-a}{A78BFA}{white}
\fm@curs{4eso-b}{6D28D9}{white}
\fm@curs{1batx-mi}{9B1C1C}{white}
\fm@curs{1batx-ccss}{EF6B6B}{white}
\fm@curs{2batx-mii}{1E6B47}{white}
\fm@curs{2batx-ccss}{2BB3A3}{white}

\newcommand{\fmcurs}[1]{%
  \@ifundefined{fm@color@#1}{\PackageError{fentmates}{Curs desconegut: #1}{}}{}%
  \edef\fm@tmp{\noexpand\definecolor{curs}{HTML}{\@nameuse{fm@color@#1}}%
               \noexpand\colorlet{curstext}{\@nameuse{fm@text@#1}}}\fm@tmp
  \colorlet{cursfosc}{curs!72!black}%
  \def\fm@logo{\fm@logodir/vertical-#1.pdf}%
  \def\fm@marca{\fm@logodir/marca-#1.pdf}}

% ---------- capçalera de la primera pàgina ----------
% #1 matèria (petit), #2 curs, #3 tipus de document i unitat, #4 títol
\newcommand{\capcalera}[4]{%
  \gdef\fm@cursnom{#2}%
  \par\noindent
  \begin{tcolorbox}[colback=curs,colframe=curs,boxrule=0pt,arc=2.2mm,width=3.5cm,height=2.15cm,
      valign=center,left=3mm,right=2mm,top=1mm,bottom=1mm,nobeforeafter,box align=center]
    \color{curstext}{\footnotesize #1}\par\vspace{1pt}{\fontsize{19}{21}\selectfont\bfseries #2\par}%
  \end{tcolorbox}\hspace{5mm}%
  \begin{minipage}[c]{\dimexpr\linewidth-3.5cm-5mm-2.9cm\relax}
    {\color{cursfosc}\small\bfseries\MakeUppercase{#3}}\par\vspace{2pt}
    {\color{navy}\fontsize{20}{23}\selectfont\bfseries #4\par}
  \end{minipage}\hfill
  \begin{minipage}[c]{2.6cm}\raggedleft\includegraphics[width=2.6cm]{\fm@logo}\end{minipage}\par
  \vspace{2pt}}

% ---------- barres de secció del color del curs ----------
\renewcommand{\seccio}[2]{\par\needspace{\fmneed\baselineskip}\bigskip
  \begin{tcolorbox}[colback=curs,colframe=curs,coltext=curstext,arc=2mm,boxrule=0pt,left=3mm,right=3mm,top=1.2mm,bottom=1.2mm]
  {\large\bfseries #1\quad #2}
  \end{tcolorbox}\vspace{-2pt}}

% ---------- capçalera i peu de les pàgines ----------
% primera pàgina: logos institucionals al peu
\fancypagestyle{primera}{\fancyhf{}\renewcommand{\headrulewidth}{0pt}
  \fancyfoot[L]{\raisebox{-0.32\height}{\includegraphics[height=0.78cm]{logo-gva.png}}\hspace{5mm}%
    \raisebox{-0.32\height}{\includegraphics[height=0.78cm]{logo-vermellar.png}}}
  \fancyfoot[R]{\scriptsize\color{muted}\href{https://creativecommons.org/licenses/by-nc-sa/4.0/deed.ca}{\ccbyncsa}\ \ CC BY-NC-SA 4.0 · fentmates.org\qquad\small\thepage}}
% la resta: el símbol a l'esquerra i el curs a la dreta, amb el títol al centre com abans
\fancyhead[L]{\raisebox{-0.3\height}{\includegraphics[height=0.62cm]{\fm@marca}}}
\fancyhead[R]{\small\bfseries\color{cursfosc}\fm@cursnom}
\renewcommand{\headrule}{\color{curs!45}\hrule height \headrulewidth}
\makeatother

\fmcurs{3eso}

\begin{document}
\thispagestyle{primera}

\capcalera{Matemàtiques}{3r ESO}{Fitxa d'exercicis · Unitat 6}{Sistemes d'equacions}
\vspace{10pt}
Nom i cognoms: \rule{8.4cm}{0.4pt}\hfill Grup: \rule{1.6cm}{0.4pt}\hfill Data: \rule{2.2cm}{0.4pt}
\vspace{8pt}

\begin{tcolorbox}[colback=soft,colframe=soft,arc=2mm,left=3mm,right=3mm,top=1.5mm,bottom=1.5mm]
\small Exercicis graduats per nivells de dificultat:\quad \estrelles{1}\ bàsic \qquad \estrelles{2}\ mitjà \qquad \estrelles{3}\ avançat.\\[2pt] Resol-los al quadern mostrant tots els passos i comprova sempre la solució en les dues equacions. Per al mètode gràfic, fes servir paper quadriculat.
\end{tcolorbox}

% =====================================================
\seccio{1}{Equacions lineals amb dues incògnites}

\ex{1} Quines d'aquestes equacions són lineals amb dues incògnites?\begin{tasks}(3)
  \task $3x-y=2$
  \task $x^2+y=1$
  \task $xy=4$
  \task $\dfrac{x}{2}+y=3$
  \task $2x+5=3x$
  \task $y=4x-1$
\end{tasks}

\ex{1} Comprova quins d'aquests parells són solució de $2x+3y=12$:\begin{tasks}(4)
  \task $(3,\,2)$
  \task $(0,\,4)$
  \task $(6,\,1)$
  \task $(-3,\,6)$
\end{tasks}

\ex{1} Aïlla la $y$ i completa una taula de valors amb $x=-2,\,-1,\,0,\,1,\,2$:\begin{tasks}(2)
  \task $x+y=5$
  \task $2x-y=1$
\end{tasks}\recorda{Cada parell $(x,\,y)$ de la taula és una solució de l'equació. Una equació amb dues incògnites té infinites solucions.}

\ex{2} Escriu tres solucions de cada equació:\begin{tasks}(2)
  \task $x-3y=6$
  \task $4x+2y=10$
\end{tasks}

\ex{2} Respon:\begin{tasks}(1)
  \task Troba $k$ perquè $(-1,\,3)$ siga solució de $4x+ky=8$.
  \task Troba $m$ perquè $(2,\,m)$ siga solució de $5x-2y=4$.
\end{tasks}

\ex{2} Fes una taula de valors i representa cada recta en uns eixos:\begin{tasks}(2)
  \task $x+2y=4$
  \task $3x-y=3$
\end{tasks}

\ex{3} Respon:\begin{tasks}(1)
  \task Troba la solució de $3x+2y=21$ en què la $y$ val 3 unitats més que la $x$.
  \task En quins punts talla la recta $x-2y=6$ els eixos de coordenades?
\end{tasks}


% =====================================================
\seccio{2}{Sistemes d'equacions i solucions}

\ex{1} Comprova si el parell $(2,\,-1)$ és solució de cada sistema:\begin{tasks}(2)
  \task $\begin{cases}x+y=1\\ 2x-y=5\end{cases}$
  \task $\begin{cases}3x+y=5\\ x-2y=3\end{cases}$
\end{tasks}\recorda{Un parell és solució del sistema si compleix \textbf{les dues} equacions alhora.}

\ex{1} Quin d'aquests parells és la solució del sistema $\begin{cases}x+y=3\\ 2x-y=3\end{cases}$? \quad $(1,\,2)$, \ $(2,\,1)$, \ $(3,\,0)$

\ex{2} Escriu un sistema de dues equacions que tinga per solució $x=-2,\ y=5$.

\ex{3} Troba $a$ i $b$ perquè $(3,\,-2)$ siga la solució del sistema $\begin{cases}ax+2y=5\\ 4x+by=6\end{cases}$.


% =====================================================
\seccio{3}{Mètode de substitució}

\ex{1} Resol per substitució (una incògnita ja està aïllada):\begin{tasks}(2)
  \task $\begin{cases}x=3y\\ x+y=8\end{cases}$
  \task $\begin{cases}y=2x+1\\ 3x+y=11\end{cases}$
\end{tasks}

\ex{1} Resol per substitució:\begin{tasks}(2)
  \task $\begin{cases}x+y=7\\ 2x-y=2\end{cases}$
  \task $\begin{cases}x-y=1\\ 3x+2y=13\end{cases}$
\end{tasks}\recorda{Aïlla la incògnita que té coeficient 1 o $-1$ i substitueix-la, \textbf{entre parèntesis}, en l'altra equació.}

\ex{2} Resol per substitució:\begin{tasks}(2)
  \task $\begin{cases}3x+y=1\\ 5x-2y=20\end{cases}$
  \task $\begin{cases}2x-y=-4\\ 4x+3y=2\end{cases}$
\end{tasks}

\ex{2} Resol per substitució i comprova la solució:\begin{tasks}(2)
  \task $\begin{cases}x+4y=-5\\ 3x-2y=13\end{cases}$
  \task $\begin{cases}5x-y=13\\ 2x+3y=12\end{cases}$
\end{tasks}

\ex{3} Resol per substitució (cap coeficient és 1: apareixeran fraccions):\begin{tasks}(2)
  \task $\begin{cases}2x+3y=12\\ 5x-4y=7\end{cases}$
  \task $\begin{cases}3x-2y=-1\\ 4x+5y=14\end{cases}$
\end{tasks}

\ex{3} En resoldre $\begin{cases}y=3-x\\ 2x-3y=1\end{cases}$ un alumne escriu $2x-9-3x=1$. Quin error ha comés? Resol el sistema correctament.


% =====================================================
\seccio{4}{Mètode d'igualació}

\ex{1} Resol per igualació:\begin{tasks}(2)
  \task $\begin{cases}y=x+5\\ y=3x-1\end{cases}$
  \task $\begin{cases}x=2y-3\\ x=y+1\end{cases}$
\end{tasks}\recorda{Aïlla \textbf{la mateixa} incògnita en les dues equacions i iguala les expressions.}

\ex{1} Resol per igualació:\begin{tasks}(2)
  \task $\begin{cases}x+y=6\\ x-y=2\end{cases}$
  \task $\begin{cases}y-2x=1\\ y+x=7\end{cases}$
\end{tasks}

\ex{2} Resol per igualació:\begin{tasks}(2)
  \task $\begin{cases}2x+y=8\\ 3x+y=11\end{cases}$
  \task $\begin{cases}x+3y=10\\ x-2y=0\end{cases}$
\end{tasks}

\ex{2} Resol per igualació i comprova la solució:\begin{tasks}(2)
  \task $\begin{cases}3x-y=7\\ x+2y=0\end{cases}$
  \task $\begin{cases}4x+y=-2\\ 2x-y=-4\end{cases}$
\end{tasks}

\ex{3} Resol per igualació (apareixeran denominadors):\begin{tasks}(2)
  \task $\begin{cases}2x+3y=7\\ 3x-2y=4\end{cases}$
  \task $\begin{cases}5x-2y=1\\ 3x+4y=11\end{cases}$
\end{tasks}

\ex{3} Resol el sistema $\begin{cases}y=\dfrac{x+1}{2}\\[6pt] y=\dfrac{2x-1}{3}\end{cases}$ pel mètode més adequat.


% =====================================================
\seccio{5}{Mètode de reducció}

\ex{1} Resol per reducció (els coeficients d'una incògnita ja són oposats):\begin{tasks}(2)
  \task $\begin{cases}x+y=9\\ x-y=3\end{cases}$
  \task $\begin{cases}3x+2y=13\\ x-2y=-1\end{cases}$
\end{tasks}\recorda{Si els coeficients d'una incògnita són oposats, en sumar les equacions aquesta incògnita desapareix.}

\ex{1} Resol per reducció:\begin{tasks}(2)
  \task $\begin{cases}2x+y=7\\ 2x-3y=-5\end{cases}$
  \task $\begin{cases}5x+3y=11\\ 2x-3y=-4\end{cases}$
\end{tasks}

\ex{2} Resol per reducció (cal multiplicar una de les equacions):\begin{tasks}(2)
  \task $\begin{cases}x+2y=4\\ 3x-y=5\end{cases}$
  \task $\begin{cases}4x+3y=10\\ 2x+5y=12\end{cases}$
\end{tasks}

\ex{2} Resol per reducció (cal multiplicar les dues equacions):\begin{tasks}(2)
  \task $\begin{cases}3x+4y=18\\ 5x-2y=4\end{cases}$
  \task $\begin{cases}2x-5y=-1\\ 3x+2y=8\end{cases}$
\end{tasks}

\ex{3} Resol per reducció:\begin{tasks}(2)
  \task $\begin{cases}4x+3y=-1\\ 5x+2y=4\end{cases}$
  \task $\begin{cases}6x-5y=7\\ 4x-3y=5\end{cases}$
\end{tasks}

\ex{3} Resol el sistema $\begin{cases}2x+3y=4\\ 3x-2y=1\end{cases}$ per \textbf{reducció doble}: elimina primer la $y$ per a trobar la $x$ i després elimina la $x$ per a trobar la $y$. Per què és útil ací?


% =====================================================
\seccio{6}{Mètode gràfic}

\ex{1} Llig la solució de cada sistema a partir del punt de tall de les dues rectes:\begin{tasks}(2)
  \task \begin{tikzpicture}[scale=0.4,baseline=(current bounding box.center)]
\draw[muted!25,very thin] (-5,-5) grid (5,5);
\draw[-latex,muted] (-5.4,0)--(5.6,0) node[right,font=\scriptsize]{$x$};
\draw[-latex,muted] (0,-5.4)--(0,5.6) node[above,font=\scriptsize]{$y$};
\node[font=\tiny,muted,below] at (-4,0) {-4}; \node[font=\tiny,muted,left] at (0,-4) {-4};
\node[font=\tiny,muted,below] at (-2,0) {-2}; \node[font=\tiny,muted,left] at (0,-2) {-2};
\node[font=\tiny,muted,below] at (2,0) {2}; \node[font=\tiny,muted,left] at (0,2) {2};
\node[font=\tiny,muted,below] at (4,0) {4}; \node[font=\tiny,muted,left] at (0,4) {4};
\begin{scope}\clip (-5,-5) rectangle (5,5);
\draw[accent,line width=1.1pt] (-6.000,10.000)--(6.000,-2.000);
\draw[teal,line width=1.1pt] (-6.000,-4.000)--(6.000,8.000);
\end{scope}
\end{tikzpicture}
  \task \begin{tikzpicture}[scale=0.4,baseline=(current bounding box.center)]
\draw[muted!25,very thin] (-5,-5) grid (5,5);
\draw[-latex,muted] (-5.4,0)--(5.6,0) node[right,font=\scriptsize]{$x$};
\draw[-latex,muted] (0,-5.4)--(0,5.6) node[above,font=\scriptsize]{$y$};
\node[font=\tiny,muted,below] at (-4,0) {-4}; \node[font=\tiny,muted,left] at (0,-4) {-4};
\node[font=\tiny,muted,below] at (-2,0) {-2}; \node[font=\tiny,muted,left] at (0,-2) {-2};
\node[font=\tiny,muted,below] at (2,0) {2}; \node[font=\tiny,muted,left] at (0,2) {2};
\node[font=\tiny,muted,below] at (4,0) {4}; \node[font=\tiny,muted,left] at (0,4) {4};
\begin{scope}\clip (-5,-5) rectangle (5,5);
\draw[accent,line width=1.1pt] (-6.000,-15.000)--(6.000,9.000);
\draw[teal,line width=1.1pt] (-6.000,2.500)--(6.000,-3.500);
\end{scope}
\end{tikzpicture}
\end{tasks}

\ex{1} Resol gràficament:\begin{tasks}(2)
  \task $\begin{cases}x+y=6\\ x-y=2\end{cases}$
  \task $\begin{cases}y=x+1\\ y=-x+5\end{cases}$
\end{tasks}\recorda{Fes una taula de valors per a cada equació (dos o tres punts), dibuixa les dues rectes en els mateixos eixos i busca el punt de tall.}

\ex{2} Resol gràficament i comprova la solució:\begin{tasks}(2)
  \task $\begin{cases}2x+y=4\\ x-y=-1\end{cases}$
  \task $\begin{cases}x+2y=0\\ 3x-y=7\end{cases}$
\end{tasks}

\ex{2} Resol gràficament:\begin{tasks}(2)
  \task $\begin{cases}y=2x-4\\ x+y=-1\end{cases}$
  \task $\begin{cases}x-3y=3\\ 2x+y=-8\end{cases}$
\end{tasks}

\ex{3} Representa aquests sistemes. Què hi observes? Quantes solucions tenen?\begin{tasks}(2)
  \task $\begin{cases}x+y=2\\ 2x+2y=6\end{cases}$
  \task $\begin{cases}x-2y=2\\ 3x-6y=6\end{cases}$
\end{tasks}


% =====================================================
\seccio{7}{Classificació dels sistemes}

\ex{1} Classifica aquests sistemes sense resoldre'ls (SCD, SCI o SI):\begin{tasks}(2)
  \task $\begin{cases}x+3y=2\\ 2x+6y=4\end{cases}$
  \task $\begin{cases}2x-y=1\\ 4x-2y=3\end{cases}$
  \task $\begin{cases}x+y=5\\ 2x-y=1\end{cases}$
  \task $\begin{cases}3x-6y=9\\ -x+2y=-3\end{cases}$
\end{tasks}\recorda{$\dfrac{a}{a'}\neq\dfrac{b}{b'}$: compatible determinat (secants). \ $\dfrac{a}{a'}=\dfrac{b}{b'}=\dfrac{c}{c'}$: compatible indeterminat (coincidents). \ $\dfrac{a}{a'}=\dfrac{b}{b'}\neq\dfrac{c}{c'}$: incompatible (paral·leles).}

\ex{1} Completa: si les rectes d'un sistema\ldots\begin{tasks}(1)
  \task es tallen en un punt, el sistema té \rule{2.5cm}{0.4pt} i és \rule{3.5cm}{0.4pt}.
  \task són paral·leles, el sistema té \rule{2.5cm}{0.4pt} i és \rule{3.5cm}{0.4pt}.
  \task coincideixen, el sistema té \rule{2.5cm}{0.4pt} i és \rule{3.5cm}{0.4pt}.
\end{tasks}

\ex{2} Resol per reducció i classifica:\begin{tasks}(2)
  \task $\begin{cases}x+y=3\\ 2x+2y=5\end{cases}$
  \task $\begin{cases}x-2y=1\\ 3x-6y=3\end{cases}$
\end{tasks}

\ex{2} Resol i classifica:\begin{tasks}(3)
  \task $\begin{cases}2x+3y=5\\ 4x+6y=10\end{cases}$
  \task $\begin{cases}x-y=2\\ -2x+2y=4\end{cases}$
  \task $\begin{cases}3x+y=2\\ x-y=2\end{cases}$
\end{tasks}

\ex{2} Escriu una equació que forme amb $2x-y=3$ un sistema:\begin{tasks}(3)
  \task compatible determinat
  \task compatible indeterminat
  \task incompatible
\end{tasks}

\ex{3} Per a quin valor de $m$ el sistema $\begin{cases}mx+2y=3\\ 3x+y=1\end{cases}$ és incompatible? Com és per als altres valors de $m$?

\ex{3} Troba $a$ i $b$ perquè el sistema $\begin{cases}2x+ay=b\\ 4x-6y=10\end{cases}$ tinga infinites solucions.


% =====================================================
\seccio{8}{Sistemes amb parèntesis i denominadors}

\ex{1} Escriu cada equació en la forma $ax+by=c$ i resol:\begin{tasks}(2)
  \task $\begin{cases}2(x-1)+y=7\\ x-(y-3)=3\end{cases}$
  \task $\begin{cases}3(x+y)-2x=10\\ 2(x-y)+y=-1\end{cases}$
\end{tasks}\recorda{Primer traiem parèntesis i denominadors, transposem i reduïm. Després triem el mètode més adequat.}

\ex{1} Traiem denominadors i resol:\begin{tasks}(2)
  \task $\begin{cases}\dfrac{x}{2}+y=4\\ x-\dfrac{y}{3}=1\end{cases}$
  \task $\begin{cases}\dfrac{x}{3}+\dfrac{y}{2}=2\\ \dfrac{x}{2}-\dfrac{y}{4}=1\end{cases}$
\end{tasks}

\ex{2} Resol:\begin{tasks}(2)
  \task $\begin{cases}\dfrac{x+1}{2}+\dfrac{y-1}{3}=3\\ x-y=-1\end{cases}$
  \task $\begin{cases}\dfrac{x-2}{3}=\dfrac{y+1}{4}\\ 2x-y=7\end{cases}$
\end{tasks}

\ex{2} Resol (multiplica per 10 per a eliminar els decimals):\begin{tasks}(1)
  \task $\begin{cases}0{,}5x+0{,}2y=2\\ x-y=-3\end{cases}$
\end{tasks}

\ex{3} Resol:\begin{tasks}(1)
  \task $\begin{cases}\dfrac{2x-1}{3}-\dfrac{y+2}{2}=-1\\ 3(x-y)-2(x+1)=-6\end{cases}$
\end{tasks}

\ex{3} Quin mètode és el més adequat per a cada sistema? Justifica-ho i resol-los.\begin{tasks}(3)
  \task $\begin{cases}y=3x-2\\ y=x+4\end{cases}$
  \task $\begin{cases}x+5y=9\\ 2x-3y=5\end{cases}$
  \task $\begin{cases}4x+3y=6\\ 6x-3y=24\end{cases}$
\end{tasks}


% =====================================================
\seccio{9}{Problemes amb sistemes d'equacions}

\ex{1} La suma de dos nombres és 56 i la seua diferència és 12. Quins són?\recorda{Tria les dues incògnites i escriu què representen, busca dues relacions per a plantejar el sistema, resol-lo i comprova que la solució té sentit en el problema.}

\ex{1} Un nombre és el triple d'un altre i entre els dos sumen 48. Quins són?

\ex{1} En un corral hi ha gallines i conills. En total hi ha 40 caps i 112 potes. Quants animals hi ha de cada tipus?

\ex{1} 3 kg de pomes i 2 kg de peres costen 8 €, i 1 kg de pomes i 3 kg de peres costen 5 €. Quant costa el quilo de cada fruita?

\ex{2} En una guardiola hi ha 22 monedes, unes de 2 € i altres de 50 cèntims, que sumen 26 €. Quantes n'hi ha de cada tipus?

\ex{2} L'entrada d'un cine costa 7 € per als adults i 4 € per als xiquets. Un dia es van vendre 85 entrades i es van recaptar 445 €. Quantes entrades de cada tipus es van vendre?

\ex{2} L'edat d'una mare és el quàdruple de la del seu fill. D'ací a 5 anys només serà el triple. Quants anys tenen ara?

\ex{2} El perímetre d'un rectangle fa 46 m i la base fa 7 m més que l'altura. Quines dimensions té?

\ex{2} Un test té 30 preguntes i les has contestades totes. Cada encert val 2 punts i cada errada en resta 1. Si has tret 39 punts, quantes n'has encertat?

\ex{2} En una classe hi ha 28 alumnes. Si arribaren 2 xiques més, hi hauria tantes xiques com xics. Quants xics i quantes xiques hi ha?

\ex{3} Volem 60 kg de cafè de 7 €/kg barrejant cafè de 9 €/kg amb cafè de 6 €/kg. Quants quilos de cada tipus hi hem de posar?

\ex{3} Les xifres d'un nombre de dues xifres sumen 11. Si les intercanviem, el nombre que resulta és 27 unitats menor. Quin és el nombre?

\ex{3} Dues ciutats estan a 75 km. Dos ciclistes ixen alhora, cadascun d'una ciutat, i van l'un cap a l'altre. Un va 5 km/h més ràpid que l'altre i es troben al cap de 3 hores. A quina velocitat va cadascun?

\ex{3} Si sumem 1 al numerador i al denominador d'una fracció, obtenim $\tfrac23$; si restem 1 a tots dos, obtenim $\tfrac12$. Quina és la fracció?

\ex{3} Un hotel té 50 habitacions, unes dobles i altres individuals, i en total 87 llits. Quantes habitacions hi ha de cada tipus?

\ex{3} El perímetre d'un triangle isòsceles fa 40 cm i cadascun dels costats iguals fa 5 cm més que la base. Quant fan els costats?

\vfill
\begin{center}\color{muted}\small\faCheckCircle[regular]\ \ Ànim! Si et costa algun exercici, repassa els apunts interactius de la unitat.\end{center}

\end{document}
